\documentclass[../main.tex]{subfiles}

\begin{document}

\ifSubfilesClassLoaded{

    %\tableofcontents
    
    %\setcounter{chapter}{}
}{}

\chapter{ Orientations }
% 代靖涵 25120222201319

\section{Orientations of Manifolds}

\begin{definition}{}{}
    Let \(  M  \) be a smooth \(  n  \)-manifold with or without boundary, endowed with a pointwise orientations.
    \begin{itemize}
        \item  If \(  \left( E_{i} \right)   \) is a local frame of \(  TM  \), we say that \(  \left( E_{i} \right)   \) is \dfntxt{(positively oriented)} if \(  \left( \left. E_1 \right|_{p},\cdots ,\left. E_{n} \right|_{p} \right)   \) is positively oriented basis for \(  T_{p}M  \) at each point \(  p\in U  \).A \dfntxt{negatively oriented} frame is defined analogously.
        \item A pointwise orientation is said to be \dfntxt{continuous} if every point of \(  M  \) is in the domain of an oriented local frame.
        \item An \dfntxt{orientation of \(  M  \)  } is a continuous pointwise orientation.
        \item We say that \(  M  \) is \dfntxt{orientable} if there exists an orientation for it, and \dfntxt{nonorientable} if not.
        \item An \dfntxt{oriented manifold} is an ordered pair pra\(  M,\mathcal{O}  \); an \dfntxt{oriented manifold with boundary is defined similarly}.   
    \end{itemize}
            
\end{definition}
\end{document}